% !TEX root = ../main.tex
\chapter{Topology}\label{ch:c01}
\begin{paracol}{2}

\begin{sources}
\begin{tabularx}{\linewidth}{@{}l l L@{}}
\toprule
\textbf{Source} & \textbf{Length} & \textbf{Content}\\
\midrule
\seg{1}{1}{Lecture 1: Topology} & 1:17:06 & Opening: the definition of spacetime as the goal of the school; why continuity needs a
structure on sets; topologies (axioms, examples, chaotic and discrete topology, standard topology on $\R^d$); open and closed
sets; continuous maps via pre-images, homeomorphisms; composition of continuous maps; the subset topology and restriction of maps.\\
\seg{1}{2}{Tutorial 1: Topology} & 0:32:35 & Worked solutions of the sheet by a tutor (not Schuller): multiple-choice questions on
topological spaces, small topologies on a four-element set, continuity of the identity, open sets in $\R$ and $\R^2$.\\
\bottomrule
\end{tabularx}

\smallskip
\textbf{Transcripts} \file{materials/transcripts/C01-V01.txt} and \file{C01-V02.txt} (YouTube captions).
\textbf{Sheet 1:} see \url{https://tales.mbivert.com/on-heraeus-winter-school-tutorials/}.
\end{sources}

\begin{reading}[Where this chapter sits in the other notes]
The same material, at a slower pace and with many more examples, is in the \emph{Geometric Anatomy of Theoretical Physics}
notes (Schuller, 2013/2015 course): the chapters ``Topological spaces: construction and purpose'' (chapter 4, where topologies,
continuity, homeomorphisms and the subset topology are treated, and where the chaotic and discrete topologies and the standard
topology reappear) and ``Topological spaces: some heavily used invariants'' (chapter 5). Lecture 2 of this course
continues with manifolds (chapter 6 there). The set theory that is used silently here (power set, pre-images) is in chapters 2
and 3 there.
\end{reading}

% ==================================================================
\section{Opening: what the school is about}\label{sec:c01-intro}

The first lecture is the central lecture of the Heraeus International Winter School on gravity and light, held in Linz
\ts{1}{1}{0:00}\ts{1}{1}{0:20}. Its topic is the structure of space and time and, from a physical point of view, the structure
of gravity theory: relativistic gravity is general relativity, and there is also the theory of relativistic matter
\ts{1}{1}{0:20}\ts{1}{1}{0:41}. The Einstein equations, which celebrate their 100th birthday in 2015, connect the matter content
of the universe with gravity \ts{1}{1}{0:41}\ts{1}{1}{1:01}. Matter of course gravitates; what is new in general relativity is
that the gravitational effect of matter is encoded in a change of the structure of spacetime itself, namely in its
\emph{curvature} \ts{1}{1}{1:01}\ts{1}{1}{1:22}. One side of the Einstein equations talks about matter, the other about gravity
\ts{1}{1}{1:22}.

\begin{history}[The Einstein equations]
\emph{Checked.} Einstein presented the final form of the field equations to the Prussian Academy of Sciences on
25~November 1915; hence the ``100th birthday'' in 2015. (Not said in class.)
\end{history}

Underlying both sides is a notion one needs in order to start talking about either: \emph{spacetime}
\ts{1}{1}{1:22}\ts{1}{1}{1:42}. Everyone has an intuitive notion (somehow space and time together as one big space), but it is
not precise enough \ts{1}{1}{1:42}\ts{1}{1}{2:03}. A precise notion, the physical key object that underlies all relativistic physics,
is the following one, written on the board \ts{1}{1}{2:03}.

\begin{definition}[Spacetime (the goal of the school)]\label{def:c01-spacetime}
\emph{Spacetime} is a four-dimensional \textcolor{defcol}{\textbf{topological manifold}} that carries a
\textcolor{defcol}{\textbf{smooth atlas}}, a \textcolor{defcol}{\textbf{torsion-free connection}}
compatible with a \textcolor{defcol}{\textbf{Lorentzian metric}} and a \textcolor{defcol}{\textbf{time orientation}}, and that
satisfies the \textcolor{defcol}{\textbf{Einstein equations}} \ts{1}{1}{2:03}\ts{1}{1}{2:23}\ts{1}{1}{2:44}.
\end{definition}
The coloured terms are the ones that will be defined precisely in the lectures; the other words are just a manner of speaking
\ts{1}{1}{2:44}\ts{1}{1}{3:05}. Schuller's roadmap, written below this definition, is of 24 lectures of the central course; up to
lecture 12 the central course runs alone (the first half of the month), then the full school starts with external lecturers and
about 50 scholarship holders from 23 countries \ts{1}{1}{3:05}\ts{1}{1}{3:25}. The colour code says in which lecture each term
is clarified: ``topological'' in the first lecture this morning; ``manifold'' and ``dimension'' in the second; all other terms in
the later colour-coded lectures; and in the remaining lectures much more will be said about how the Einstein equations work, but
first the notion of spacetime has to be understood \ts{1}{1}{3:25}\ts{1}{1}{3:46}\ts{1}{1}{4:06}.

\paragraph{Organisation.} The schedule is tight; students may interrupt at any time and get very brief answers
\ts{1}{1}{4:06}\ts{1}{1}{4:27}. Every afternoon after the lectures there is a tutorial with three supervisors; the sheets are
picked up from the stacks, there is nothing to prepare, and the exercises are very simple ones that make one accustomed to the
underlying notions; one makes much more of the school by going \ts{1}{1}{4:27}\ts{1}{1}{4:47}\ts{1}{1}{5:07}.

A student immediately takes up the invitation to interrupt: is ``all modern physics'' not a rather broad claim? Yes, of course:
whenever Schuller speaks about physics he needs to be a bit grand. What he means is all modern \emph{fundamental relativistic}
physics \ts{1}{1}{5:07}\ts{1}{1}{5:27}.

% ==================================================================
\section{Topologies}\label{sec:c01-topology}

\subsection{Why we need a topology}

Lecture one is about topological spaces \ts{1}{1}{5:27}. At the coarsest level, if one does not look closely, spacetime is a
\emph{set}: it consists of points, which are the elements of the set, and that is it \ts{1}{1}{5:49}. This is not enough to talk
even about the simplest notions of classical physics, namely continuity, i.e.\ continuity of maps \ts{1}{1}{6:10}. Why do we, as
physicists, want continuity? In classical physics curves do not jump: a particle running somewhere (later: in spacetime) is not
suddenly continuing in a totally different place, so one must require, for instance, continuity of curves
\ts{1}{1}{6:32}\ts{1}{1}{6:52}. A set carries not enough structure to say whether a curve on it is continuous \ts{1}{1}{6:52}\ts{1}{1}{7:12}.

One could imagine many structures on a set that allow continuity (for instance a distance measure), but one wants to be very
minimal and economical in order not to introduce undue assumptions: the \emph{weakest} structure that has to be established
by hand on a set and allows a good definition of continuity of maps \ts{1}{1}{7:12}\ts{1}{1}{7:33}\ts{1}{1}{7:53}. Mathematicians
know that this weakest structure is a so-called \emph{topology}; this is the reason why the classical physicist looks at
topologies \ts{1}{1}{7:53}\ts{1}{1}{8:14}. All of this is the motivation, ``the little cloud'' one may keep in mind; from now on
the statements are precise \ts{1}{1}{8:14}\ts{1}{1}{8:36}.

\subsection{Definition and the three axioms}

\begin{definition}[New concepts: topology]\label{def:c01-topology}
Let $M$ be a set. A \textbf{topology} $\mathcal O$ on $M$ is a subset $\mathcal O\subseteq\mathcal P(M)$ of the \emph{power set}
$\mathcal P(M)$ (the set of all subsets of $M$), so a set of sets, that satisfies three axioms
\ts{1}{1}{8:36}\ts{1}{1}{8:57}\ts{1}{1}{9:18}\ts{1}{1}{9:39}:
\begin{enumerate}
  \item $\emptyset\in\mathcal O$ and $M\in\mathcal O$ \ts{1}{1}{10:02}\ts{1}{1}{10:22};
  \item $U,V\in\mathcal O\ \Rightarrow\ U\cap V\in\mathcal O$ \ts{1}{1}{10:42}\ts{1}{1}{11:04};
  \item for an arbitrary index set $A$ (finite, countable, or as large as $\R$ or more):
        $\{U_\alpha\}_{\alpha\in A}\subseteq\mathcal O\ \Rightarrow\ \bigcup_{\alpha\in A}U_\alpha\in\mathcal O$
        \ts{1}{1}{11:25}\ts{1}{1}{11:45}\ts{1}{1}{12:06}\ts{1}{1}{12:26}.
\end{enumerate}
\end{definition}
Here $\emptyset$ is a subset of $M$ and $M$ is an (improper) subset of $M$; the empty set and $M$ must always be chosen to belong to
$\mathcal O$, otherwise $\mathcal O$ does not qualify as a topology for $M$. The same holds if some intersection of two chosen
sets has not been chosen \ts{1}{1}{10:02}\ts{1}{1}{10:22}\ts{1}{1}{11:04}. The third axiom looks deceptively similar to the
second one, but with a union instead of an intersection and an arbitrary instead of a pair-wise collection; it is
``somewhat different''. Each time one has to \emph{check} the three axioms \ts{1}{1}{9:39}\ts{1}{1}{11:25}\ts{1}{1}{12:26}.
This is the key definition of topology; the questions are what it is all about, and whether we can have an example
\ts{1}{1}{12:46}. (The letter $\mathcal O$ will be explained later: it stands for \emph{open}, see below.)

\subsection{Examples on a three-element set}

\begin{example}[A topology on $\{1,2,3\}$]\label{ex:c01-O1}
Let $M=\{1,2,3\}$ and $\mathcal O_1=\{\emptyset,\{1,2,3\}\}$; these are the only subsets of $M$ that $\mathcal O_1$ contains
\ts{1}{1}{13:09}\ts{1}{1}{13:31}. It is a topology \ts{1}{1}{13:52}. The empty set and $M$ lie in it. For the second axiom take
any two members as $U$ and $V$: $\emptyset\cap\emptyset=\emptyset$, $\emptyset\cap M=\emptyset$, $M\cap\emptyset=\emptyset$ (the
intersection is symmetric) and $M\cap M=M$, all already in $\mathcal O_1$. For the third axiom the same holds: all unions of
members are again $\emptyset$ or $M$ \ts{1}{1}{14:13}\ts{1}{1}{14:33}.
\end{example}
\begin{example}[Not a topology]\label{ex:c01-O2}
Same set, $\mathcal O_2=\{\emptyset,\{1\},\{2\},\{1,2,3\}\}$ \ts{1}{1}{15:15}. It is \emph{not} a topology: the third axiom fails, since
$\{1\}\cup\{2\}=\{1,2\}\notin\mathcal O_2$ \ts{1}{1}{15:35}. (A student's variant of the example is dealt with in the same way:
$\{1,2\}$ is a subset of $M$, but it is not an element of the collection; the union is not in it \ts{1}{1}{17:41}\ts{1}{1}{18:01}\ts{1}{1}{18:23}.)
\end{example}

A student asks whether a topology has to be countable. No, in fact it won't be, and not even if $M$ is countable: if
$M=\N$, the power set of $M$ is of the size of the reals, and one may choose $\mathcal O$ to be the entire power set (the next
example) \ts{1}{1}{14:55}\ts{1}{1}{15:15}.

This is a very simple game as long as one can write the topology down explicitly \ts{1}{1}{15:57}.

\subsection{The two extreme topologies}

\begin{definition}[New concepts: chaotic and discrete topology]\label{def:c01-extreme}
For \emph{any} set $M$ (the integers, the reals, whatever) there are two topologies that can always be written down
\ts{1}{1}{15:57}\ts{1}{1}{16:18}:
\begin{align*}
  \mathcal O_{\mathrm{chaotic}}&=\{\emptyset,M\} && \text{the \textbf{chaotic topology}} \ts{1}{1}{16:18}\ts{1}{1}{16:38},\\
  \mathcal O_{\mathrm{discrete}}&=\mathcal P(M) && \text{the \textbf{discrete topology}} \ts{1}{1}{16:38}\ts{1}{1}{16:59}.
\end{align*}
\end{definition}
That the chaotic topology is a topology can be checked for any $M$ (as in \cref{ex:c01-O1}); the discrete one is a topology
because the intersection and the union of two subsets are again subsets, and all subsets are in it \ts{1}{1}{16:18}\ts{1}{1}{16:59}.
Why bother with them? Because they are both utterly useless, but they are \emph{extreme} cases: the chaotic topology has the
fewest elements possible, the discrete one the maximum. They serve as test cases, as extreme examples, to see on what a notion
in topology depends \ts{1}{1}{16:59}\ts{1}{1}{17:20}\ts{1}{1}{17:41}.

\subsection{The standard topology on \texorpdfstring{$\R^d$}{Rd}}

There is a very important example, used throughout, that reconciles our undergraduate intuition of continuity (to be defined
soon) with the notion of topology: the \emph{standard topology} \ts{1}{1}{18:23}\ts{1}{1}{18:43}. Whereas the chaotic and
discrete topologies exist on any set, this one can only be defined on $\R^d=\R\times\dots\times\R$, the set of all $d$-tuples
$(p_1,\dots,p_d)$ of real numbers \ts{1}{1}{18:43}\ts{1}{1}{19:03}\ts{1}{1}{19:24}. One could put the chaotic or the discrete
topology on $\R^d$, or a great number of others; this is the one that is usually chosen without saying so
\ts{1}{1}{19:24}\ts{1}{1}{19:45}. It contains uncountably many sets, which cannot be written down explicitly: it is defined
implicitly in two steps, and then one checks that this defines a topology uniquely \ts{1}{1}{20:06}\ts{1}{1}{20:26}\ts{1}{1}{20:47}.

\begin{definition}[New concepts: soft balls and the standard topology]\label{def:c01-standard}
\textbf{Step $\alpha$.} For $p\in\R^d$ and $r\in\R^+$ the \textbf{soft ball} (open ball) of radius $r$ and centre $p$ is the set
\[
  B_r(p):=\Bigl\{(q_1,\dots,q_d)\in\R^d\ \Big|\ \sum_{i=1}^{d}(q_i-p_i)^2<r^2\Bigr\}
\]
\ts{1}{1}{20:47}\ts{1}{1}{21:08}\ts{1}{1}{21:28}\ts{1}{1}{21:48}.\\
\textbf{Step $\beta$.} A set $U\subseteq\R^d$ lies in the \textbf{standard topology} $\mathcal O_{\mathrm{standard}}$ \emph{by definition}
if and only if for every point $p\in U$ there exists an $r\in\R^+$ such that $B_r(p)\subseteq U$
\ts{1}{1}{22:09}\ts{1}{1}{22:29}\ts{1}{1}{22:50}\ts{1}{1}{23:11}.
\end{definition}
Remark: what is written in step $\alpha$ looks like the Euclidean norm, but no norm is spoken of and none is needed; the
expression is just defined on tuples, and no vector space structure on $\R^d$ is required (to write down a norm one would need
one) \ts{1}{1}{21:48}\ts{1}{1}{22:09}.

\paragraph{Picture.} Think of a subset of $\R^2$ whose boundary does not belong to it (roughly speaking). Is it in the standard
topology? Yes: whatever point of $U$ one picks, one finds a radius small enough such that the soft ball (strict inequality, ``not an
equal sign'') lies entirely in $U$. If instead part of the boundary still belonged to the set, one could pick a point on it, and
however small but positive the radius, the ball would not lie fully in the set. So that set is \emph{not} in the topology
\ts{1}{1}{23:11}\ts{1}{1}{23:34}\ts{1}{1}{23:55}\ts{1}{1}{24:15}\ts{1}{1}{24:36}. These are the pictures one would have called
open and non-open sets: precisely the reason for defining this special topology this way is to do on $\R^d$ what one always knew
\ts{1}{1}{25:17}\ts{1}{1}{25:37}.

Showing that the implicitly defined collection satisfies the three axioms is a simple, few-line exercise that Schuller leaves out
for reasons of time (one might think about it tonight) \ts{1}{1}{24:36}\ts{1}{1}{24:57}\ts{1}{1}{25:17}. Here is the argument, which was
\emph{not said in class}:
\begin{proposition}[Standard topology is a topology]\label{prop:c01-std-ok}
$\mathcal O_{\mathrm{standard}}$ satisfies the three axioms of \cref{def:c01-topology}.
\end{proposition}
\begin{proof}[Proof (not said in class)]
$\emptyset$: the condition ``for every $p\in\emptyset$'' is empty, so $\emptyset\in\mathcal O_{\mathrm{standard}}$; $\R^d$: any ball
lies in $\R^d$. Intersection: for $p\in U\cap V$ there are $r_1,r_2$ with $B_{r_1}(p)\subseteq U$, $B_{r_2}(p)\subseteq V$, and
$B_{\min(r_1,r_2)}(p)\subseteq U\cap V$. Union: if $p\in\bigcup_\alpha U_\alpha$ then $p\in U_\beta$ for some $\beta$, so some
ball around $p$ lies in $U_\beta\subseteq\bigcup_\alpha U_\alpha$. (No triangle inequality is even needed.)
\end{proof}

\begin{coursenote}[Why such a long motivation]
Schuller adds that one should not be deterred by the abstractness: for a mathematician this is just how things go; a physicist has
to think about where these assumptions are usually made implicitly. One task in order to understand spacetime is to make the
implicit assumptions \emph{explicit}; Einstein did not understand many of them when he wrote down general relativity, and it took
him many years to overcome intuitive, even ``obvious'', notions that nevertheless do not apply \ts{1}{1}{26:40}\ts{1}{1}{27:00}\ts{1}{1}{27:20}.
\end{coursenote}

Later $M$ will be the set of points of spacetime, and a big question will be \emph{which topology should the physical spacetime,
or rather our mathematical model of it, carry?} There is a huge choice: anything that fits into the scheme of the three axioms is a
topology, and any such one suffices to define continuity of maps; recall that this is why we defined topologies
\ts{1}{1}{25:37}\ts{1}{1}{25:58}\ts{1}{1}{26:19}. Schuller asks to be trusted that none of the mathematical notions introduced is
unnecessary for the physics: later he will be able to point at each notion and say what the physical effect of this or that choice
is. It was in the second half of the twentieth century that Wald, Ehlers and others recognised that general relativity (spacetime,
gravity and matter) must be studied on the differential-geometry side, with the rigour with which mathematicians study differential
geometry, to come to terms with its subtle notions; the language is self-contained, ``the cleanest way to introduce it''
\ts{1}{1}{27:20}\ts{1}{1}{27:41}\ts{1}{1}{28:03}\ts{1}{1}{28:23}.

\subsection{Terminology: topological space, open and closed sets}

A topology is something one \emph{chooses}; it is not automatically there, and whenever one chooses, one provides structure
\ts{1}{1}{28:44}. What follows is terminology, not definitions \ts{1}{1}{28:44}\ts{1}{1}{29:07}.

\begin{definition}[New concepts: topological space, open set, closed set]\label{def:c01-open-closed}
\begin{itemize}
  \item A \textbf{topological space} is a pair $(M,\mathcal O)$ of a set $M$ and a topology $\mathcal O$ on it: more than a set, a set with
        additional information \ts{1}{1}{29:27}\ts{1}{1}{29:47}.
  \item A subset $U\subseteq M$ with $U\in\mathcal O$ is called an \textbf{open set} \ts{1}{1}{30:07}.
  \item A subset $A\subseteq M$ is called \textbf{closed} if its complement $M\setminus A$ lies in $\mathcal O$ \ts{1}{1}{30:51}\ts{1}{1}{31:12}.
\end{itemize}
\end{definition}
So the topology is the set of open sets; but it is not that one knows what ``open'' is and then has the topology: rather, ``open'' is
\emph{defined} by having chosen a topology \ts{1}{1}{30:07}\ts{1}{1}{30:28}. About the word ``set'': one can make it a problem
what a set is (a collection of elements, but what is an element, what is a collection?); one would have to write down axioms of
set theory; this is not done, one keeps the intuitive notion, and in doubt the reference is Zermelo--Fraenkel set theory with the
axiom of choice \ts{1}{1}{29:07}\ts{1}{1}{29:27}.

\begin{remark}[Warning: open and closed are not opposites]\label{rem:c01-warning}
One might think that ``closed'' is the opposite of ``open'', but not in any way \ts{1}{1}{31:12}\ts{1}{1}{31:32}. There are sets that are
open and closed at the same time: the empty set is in $\mathcal O$, and $M\setminus\emptyset=M$ is in $\mathcal O$ too, so $\emptyset$ is open
and closed. There are sets that are open but not closed, sets that are closed but not open, and sets that are neither
\ts{1}{1}{31:32}\ts{1}{1}{31:53}. In plain English, a priori the two notions have nothing to do with each other; only a
chosen topology connects them through the definition \ts{1}{1}{31:53}\ts{1}{1}{32:14}. Schuller puts up a ``warning sign'' here and will do so
several times in the lectures: from open it does \emph{not} follow that a set is closed, neither the other way round, despite the
deceptive terminology \ts{1}{1}{32:14}.
\end{remark}
\begin{aside}[Vacuum physicists]
Whether the everyday meaning of the words is useful depends on what one looks at: for instance, vacuum physicists find examples
where it is. Nevertheless the extreme cases belong to the whole picture, one cannot understand the other bits if one does not
understand them \ts{1}{1}{32:14}\ts{1}{1}{32:42}.
\end{aside}

% ==================================================================
\section{Continuous maps}\label{sec:c01-continuity}

\subsection{Why did we buy a topology?}

Having a topology gives a notion of \emph{continuous maps}; that is what it does us as a service \ts{1}{1}{33:03}\ts{1}{1}{33:24}.
First recall maps: a map $f\colon M\to N$ goes from a set $M$, the \emph{domain}, to a set $N$, the \emph{target}, taking every point of
the domain to some point of the target \ts{1}{1}{33:24}\ts{1}{1}{33:45}\ts{1}{1}{34:06}. Two points may be taken to the same point; not every
point of the target has to be hit. If points of the target are not hit, the map is not \emph{surjective}; if some target points are hit
more than once, it is not \emph{injective}; if it is both, it is \emph{bijective}. These notions need no further structure on the
sets, a map is defined between two sets only \ts{1}{1}{34:06}\ts{1}{1}{34:26}\ts{1}{1}{34:48}\ts{1}{1}{35:09}\ts{1}{1}{35:29}\ts{1}{1}{35:49}.

Given such a map one may ask: is it continuous, does it make jumps? One might have an intuition, but the answer is \emph{it depends}.
It depends \emph{by definition} (not by some divine intervention) on which topologies are chosen on the domain and on the target; to decide
continuity one has to decide on a topology on $M$ and on one on $N$ \ts{1}{1}{35:49}\ts{1}{1}{36:10}\ts{1}{1}{36:30}\ts{1}{1}{36:51}\ts{1}{1}{37:12}.

\subsection{Definition via pre-images}

Notation: for a map one writes $f\colon M\to N$ (name, domain, target) and, for a specific element $m\in M$, $m\mapsto f(m)$, the \emph{image}
of $m$ in $N$, e.g.\ $f(m)=m^2$ for real numbers; in the worst case one has to draw a picture, but a definite prescription is always
assumed to be given \ts{1}{1}{37:32}\ts{1}{1}{38:13}\ts{1}{1}{38:33}\ts{1}{1}{38:54}\ts{1}{1}{39:15}.

\begin{definition}[New concepts: pre-image and continuous map]\label{def:c01-continuous}
Let $f\colon M\to N$ be a map of sets. The \textbf{pre-image} of a subset $V\subseteq N$ is the subset of the domain
\[
  \operatorname{preim}_f(V)\;=\;f^{-1}(V)\;:=\;\{m\in M\mid f(m)\in V\}
\]
\ts{1}{1}{41:01}\ts{1}{1}{41:22}\ts{1}{1}{41:43}. Now let $(M,\mathcal O_M)$ and $(N,\mathcal O_N)$ be topological spaces (possibly with
totally different topologies; if the sets differ, the topologies can never be the same). The map $f$ is called
\textbf{continuous} (more fully: continuous with respect to $\mathcal O_M$ and $\mathcal O_N$) if
\[
  \forall V\in\mathcal O_N:\quad \operatorname{preim}_f(V)\in\mathcal O_M
\]
\ts{1}{1}{37:32}\ts{1}{1}{37:53}\ts{1}{1}{39:36}\ts{1}{1}{39:58}\ts{1}{1}{40:19}\ts{1}{1}{40:40}.
\end{definition}
This must be the simplest definition ever: no epsilons, no deltas \ts{1}{1}{43:47}\ts{1}{1}{44:08}. The pre-image is defined for any subset,
open or not; the definition of continuity \emph{requires} that one takes an open subset in the target and asks whether its pre-image,
a subset of the domain, is open in the domain \ts{1}{1}{41:01}\ts{1}{1}{41:22}. The mnemonic phrase: \textbf{a map is continuous if
the pre-images of open sets are open sets}; the first ``open sets'' refer to the topology on the target, the second to the topology on the domain
\ts{1}{1}{47:12}\ts{1}{1}{47:32}\ts{1}{1}{47:53}\ts{1}{1}{48:13}\ts{1}{1}{48:33}.

\paragraph{Pre-image is not the inverse.} Using a picture of a map between two small sets: if $V$ consists of two points of the target, one follows the
arrows backwards and the pre-image is made of the points of the domain that land in $V$; the point that does not land in $V$ is not
included \ts{1}{1}{41:43}\ts{1}{1}{42:04}\ts{1}{1}{42:25}\ts{1}{1}{42:45}. Careful, it is not ``kind of the inverse'': take a target point
that is hit twice; its pre-image is both points of the domain that are mapped to it \ts{1}{1}{42:45}\ts{1}{1}{43:06}. So even where the map is not
invertible, the pre-image makes sense (for a single point of the invertible part of the map it would be a single point, but not here), and the pre-image
is a subset of the domain \ts{1}{1}{43:06}\ts{1}{1}{43:26}\ts{1}{1}{43:47}.

\paragraph{Student questions on the pre-image.}
\begin{itemize}
  \item If $V$ is the set containing only the target point that is not hit, what is its pre-image? The \emph{empty set}. The empty set is open; so
        if such a one-element set happens to be open in the target, this does not destroy continuity. That is why it is so important to have the
        empty set in the topology: otherwise, with this definition, only surjective maps could be continuous, which is nonsense
        \ts{1}{1}{44:08}\ts{1}{1}{44:29}\ts{1}{1}{44:49}.
  \item In the picture, where does one see the topology? Nowhere, that is the point. The pre-image is defined at the level of maps, i.e.\ of sets; for continuity one
        chooses an open set in the target and checks whether it is open in the domain. For this one needs the topologies; for the pre-images one does not.
        Even if one understood all pre-images of all subsets, one still could not say whether the map is continuous until one has decided on topologies on both sides
        \ts{1}{1}{44:49}\ts{1}{1}{45:09}\ts{1}{1}{45:29}\ts{1}{1}{45:50}.
  \item Can a map ``make a connection'' between the topologies? In fact, if the map is bijective, the pre-image of $V$ can be written $f^{-1}(V)$ (the image under the inverse map).
        If it is continuous in \emph{both} directions, then not only pre-images but also images of open sets are open, since the pre-image under the inverse is the image under the
        map \ts{1}{1}{45:50}\ts{1}{1}{46:10}\ts{1}{1}{46:31}.
\end{itemize}

\begin{definition}[New concepts: homeomorphism]\label{def:c01-homeo}
A map $f\colon M\to N$ between topological spaces that is bijective, continuous, and whose inverse is continuous, is called a
\textbf{homeomorphism} (German: Hom\"oomorphismus; Schuller comments on the English pronunciation).
Homeomorphisms are the structure-preserving maps of topology \ts{1}{1}{46:31}\ts{1}{1}{46:51}.
\end{definition}
That is a deep point, because mathematicians always classify structures by identifying them whenever structure-preserving maps
between them exist \ts{1}{1}{46:51}\ts{1}{1}{47:12}.

\begin{coursenote}[Name]
``Homeomorphism'' comes from the Greek \emph{homoios} (similar) and \emph{morphe} (shape): same shape. (Not said in class.)
\end{coursenote}

\subsection{An example: a map, and its inverse}

\begin{example}[Continuous, but the inverse is not]\label{ex:c01-f-g}
\textbf{Part A.} Let $M=N=\{1,2\}$ as sets, and $f\colon M\to N$ the map $1\mapsto2$, $2\mapsto1$ (``why would I not do that?'').
Equip $M$ with $\mathcal O_M=\{\emptyset,\{1\},\{2\},\{1,2\}\}$ (a topology) and $N$ with the different topology
$\mathcal O_N=\{\emptyset,\{1,2\}\}$ \ts{1}{1}{50:16}\ts{1}{1}{50:42}\ts{1}{1}{51:02}\ts{1}{1}{51:23}\ts{1}{1}{51:43}. Is $f$ continuous (with respect to these
topologies)? One needs the pre-images of the open sets of the target, which can be completely enumerated: $f^{-1}(\emptyset)=\emptyset$ and
$f^{-1}(\{1,2\})=M$, and both lie in $\mathcal O_M$ by the axioms. So $f$ is continuous \ts{1}{1}{52:04}\ts{1}{1}{52:24}\ts{1}{1}{52:45}\ts{1}{1}{53:06}.

\textbf{Part B.} Now the map $g$ in the opposite direction, $g\colon N\to M$, $1\mapsto2$, $2\mapsto1$, with the topologies swapped: the domain
now carries $\mathcal O_N$ and the target $\mathcal O_M$. As a set map it is the same as $f$ (and can be viewed as $f^{-1}$), but not once the topologies
are kept \ts{1}{1}{53:27}\ts{1}{1}{53:48}\ts{1}{1}{56:13}\ts{1}{1}{56:34}\ts{1}{1}{56:55}. One might expect that the inverse of a continuous bijective map is
continuous; check: the pre-image of the open set $\{1\}$ of the target is $\{2\}$, which is not in the topology $\mathcal O_N$ of the domain. So the
inverse is \emph{not} continuous \ts{1}{1}{56:55}\ts{1}{1}{57:16}\ts{1}{1}{57:36}\ts{1}{1}{57:58}.
\end{example}
It is even worse: had $1$ and $2$ not been exchanged, one would have the identity map from a space with one topology to the same set with another topology;
it can be continuous in one direction while its inverse, again the identity, is not, because different topologies were chosen. ``Do you see how
dangerous it is not to talk explicitly about the topologies?'' \ts{1}{1}{57:58}\ts{1}{1}{58:19}.

\paragraph{The pre-image of the empty set (student question).} A student asks why the pre-image of the empty set is the empty set: because of what a map is.
The pre-image is $\{m\in M\mid f(m)\in\emptyset\}$; since $f$ is a map, every $m$ has an image in the target, which cannot be in the empty set; so no $m$
satisfies the condition. Even if the notion of map allowed some points not to be mapped, such a point would still not belong to the pre-image. In the end it is
the definition of a map \ts{1}{1}{53:06}\ts{1}{1}{54:08}\ts{1}{1}{54:29}\ts{1}{1}{54:50}\ts{1}{1}{55:10}\ts{1}{1}{55:32}\ts{1}{1}{55:52}\ts{1}{1}{56:13}.

\paragraph{The standard topology is nothing special.} A student asks whether from now on the standard topology is used. No: it can only be defined on
$\R^d$, and the sets in this example are not of that form. The name is deceptive, as if there were anything mathematically special about it; there is not,
except that one likes it because it is the first topology one met implicitly at the beginning of one's studies. That is not a good reason for spacetime to
carry it. Hence the whole discussion is kept general unless one says which topology is chosen \ts{1}{1}{48:33}\ts{1}{1}{48:54}\ts{1}{1}{49:14}\ts{1}{1}{49:35}\ts{1}{1}{49:55}\ts{1}{1}{50:16}.

\paragraph{Open versus closed once more (student question).} A student asks for the difference between open and closed in the examples. In $\mathcal O_M=\{\emptyset,\{1\},\{2\},\{1,2\}\}$
the set $\{2\}$ is open, and it is also \emph{closed}: its complement $\{1,2\}\setminus\{2\}=\{1\}$ lies in the topology. This has nothing to do with the
boundary being there or not; ``boundary'' can be defined for any topology, and the notions carry over, but our \emph{pictures} of them do not necessarily
\ts{1}{1}{58:42}\ts{1}{1}{59:03}\ts{1}{1}{59:24}\ts{1}{1}{59:44}\ts{1}{1}{1:00:05}.

\paragraph{Curves.} In the beginning a curve making a jump was drawn. To make sense of it, a curve is a map from the real line (the curve parameter) into some set,
and one must put a topology on the parameter space and on the space in which the curve runs: jumping is translated as ``not continuous''
\ts{1}{1}{1:00:05}\ts{1}{1}{1:00:26}\ts{1}{1}{1:00:46}. \textbf{Convention:} in all the lectures, whenever $\R$, $\R\times\R$ or $\R^d$ appear and nothing else is said,
the standard topology is implicitly assumed, so that it need not be written all the time. The other spaces of the lectures will not be $\R^d$ or $\R^4$:
spacetime will be a topological space all by itself \ts{1}{1}{1:00:46}\ts{1}{1}{1:01:07}\ts{1}{1}{1:01:27}\ts{1}{1}{1:01:48}.

% ==================================================================
\section{Composition of continuous maps}\label{sec:c01-composition}

Composition of maps is defined at the set level \ts{1}{1}{1:01:48}\ts{1}{1}{1:02:09}. Given $f\colon M\to N$ and $g\colon N\to P$ (the name of the map is sometimes written on the
arrow), define the map \textbf{$g$ after $f$}, $g\circ f\colon M\to P$, by $(g\circ f)(m):=g(f(m))$: first apply $f$, then $g$. Why introduce a notation for a
map when $g(f(m))$ already describes the image of a point? Because one wants to talk about \emph{the map itself}
\ts{1}{1}{1:02:09}\ts{1}{1}{1:02:31}\ts{1}{1}{1:02:52}\ts{1}{1}{1:03:12}\ts{1}{1}{1:03:33}. The question is: is the composition of continuous maps again continuous?

\begin{theorem}[Composition of continuous maps]\label{thm:c01-composition}
Let $(M,\mathcal O_M)$, $(N,\mathcal O_N)$, $(P,\mathcal O_P)$ be topological spaces. If $f\colon M\to N$ and $g\colon N\to P$ are continuous, then
$g\circ f$ is continuous \ts{1}{1}{1:03:33}\ts{1}{1}{1:03:53}.
\end{theorem}
\begin{proof}
Let $V\in\mathcal O_P$ be any open set. Then
\begin{align*}
  (g\circ f)^{-1}(V)&=\{m\in M\mid g(f(m))\in V\}\\
  &=\{m\in M\mid f(m)\in g^{-1}(V)\}\\
  &=f^{-1}\bigl(g^{-1}(V)\bigr)
\end{align*}
\ts{1}{1}{1:03:53}\ts{1}{1}{1:04:14}\ts{1}{1}{1:04:35}\ts{1}{1}{1:04:55}\ts{1}{1}{1:05:15}. Since $g$ is continuous and $V$ is open in $P$, $g^{-1}(V)\in\mathcal O_N$; since $f$ is
continuous, $f^{-1}(g^{-1}(V))\in\mathcal O_M$. So the pre-image of an arbitrary open set under $g\circ f$ is open, which is what was to be shown
\ts{1}{1}{1:05:15}\ts{1}{1}{1:05:36}.
\end{proof}
Compare with doing this with epsilons and deltas in $\R^d$: here only the definition of continuity on topological spaces is used, and the result holds always
\ts{1}{1}{1:05:36}\ts{1}{1}{1:05:56}.

% ==================================================================
\section{Inheriting a topology: the subset topology}\label{sec:c01-subset}

If you are old enough, you know that inheritance is a good thing: if one has one, two or three topological spaces (``three parents''), one may construct a new one from what one
already has. There are many useful ways to inherit topologies, each way having its advantage in a certain situation; the one particularly important for spacetime physics
(``so that the physicists keep motivated'') is the following \ts{1}{1}{1:05:56}\ts{1}{1}{1:06:21}\ts{1}{1}{1:06:43}\ts{1}{1}{1:07:04}\ts{1}{1}{1:07:25}.

Let $(M,\mathcal O_M)$ be a topological space and $S\subseteq M$ \emph{any} subset (not necessarily an open one). Can one construct on $S$ a topology from $\mathcal O_M$? One
can always invent a topology on $S$, but one wants a special one that one gets from $M$ \ts{1}{1}{1:07:25}\ts{1}{1}{1:07:46}\ts{1}{1}{1:08:07}\ts{1}{1}{1:08:28}.

\begin{definition}[New concepts: subset topology]\label{def:c01-subset}
The \textbf{subset topology} (inherited from the superset $M$), also written $\mathcal O_M|_S$, is the subset of $\mathcal P(S)$
\[
  \mathcal O|_S:=\{\,U\cap S\mid U\in\mathcal O_M\,\}
\]
\ts{1}{1}{1:08:28}\ts{1}{1}{1:08:49}\ts{1}{1}{1:09:09}\ts{1}{1}{1:09:31}.
\end{definition}
The claim is that it really is a topology, not only called one \ts{1}{1}{1:09:31}.
\begin{proposition}[The subset topology is a topology]\label{prop:c01-subset-ok}
$\mathcal O|_S$ satisfies the three axioms \ts{1}{1}{1:09:31}\ts{1}{1}{1:09:51}.
\end{proposition}
\begin{proof}
\begin{itemize}
  \item $\emptyset=\emptyset\cap S$ with $\emptyset\in\mathcal O_M$, so $\emptyset\in\mathcal O|_S$ \ts{1}{1}{1:09:51}\ts{1}{1}{1:10:12}. Also $S=M\cap S$ because $S\subseteq M$, and $M\in\mathcal O_M$,
        so $S\in\mathcal O|_S$ \ts{1}{1}{1:10:12}\ts{1}{1}{1:10:32}.
  \item Let $A,B\in\mathcal O|_S$. Then there are $\tilde A,\tilde B\in\mathcal O_M$ with $A=\tilde A\cap S$, $B=\tilde B\cap S$ \ts{1}{1}{1:10:32}\ts{1}{1}{1:10:52}\ts{1}{1}{1:11:13}. By associativity of the intersection and $S\cap S=S$,
        \[
           A\cap B=(\tilde A\cap S)\cap(\tilde B\cap S)=(\tilde A\cap\tilde B)\cap S,
        \]
        and $\tilde A\cap\tilde B\in\mathcal O_M$, so $A\cap B\in\mathcal O|_S$ \ts{1}{1}{1:11:13}\ts{1}{1}{1:11:33}.
  \item For unions the same argument works (``I am not going to write it down for time reasons''): $\bigcup_\alpha(\tilde A_\alpha\cap S)=\bigl(\bigcup_\alpha\tilde A_\alpha\bigr)\cap S$
        \ts{1}{1}{1:11:33}\ts{1}{1}{1:11:54}.
\end{itemize}
\end{proof}

\paragraph{Why this topology: restriction of maps.} Very often it is easy to decide whether a map $f\colon M\to N$, with known topologies on both sides, is continuous (for instance
with the undergraduate analysis notion); and one often wants to look at a subset $S\subseteq M$ and at the \emph{restriction} $f|_S\colon S\to N$ of the map to it
\ts{1}{1}{1:12:15}\ts{1}{1}{1:12:37}\ts{1}{1}{1:12:58}\ts{1}{1}{1:13:18}.

\begin{example}[A temperature on a hot plate, and a mouse on a wire]\label{ex:c01-mouse}
A map $T\colon\R^2\to\R$, e.g.\ a temperature distribution (isotherms) on a hot plate. For some reason one is not interested in the temperature on the whole plate, but on a little wire that
runs through it, on which a mouse is mounted with some little wheels; the mouse is only interested in the temperature where it can pass: the subset
$S\subseteq\R^2$ is the wire \ts{1}{1}{1:13:18}\ts{1}{1}{1:13:39}\ts{1}{1}{1:14:00}\ts{1}{1}{1:14:21}. Continuity of $T$ is decided by the undergraduate notion; but if the wire is
complicated, the restriction $T|_S$ ``is quite a beast'' \ts{1}{1}{1:14:21}\ts{1}{1}{1:14:42}.
\end{example}
If $S$ is given just some topology one has come up with, one cannot conclude from the continuity of $T$ that $T|_S$ is continuous; it depends on the topology chosen on $S$. But if one chooses
on $S$ the topology inherited from $\R^2$, then one is guaranteed that the restriction $T|_S\colon S\to\R$ is again continuous. It is therefore a very natural thing to consider
\ts{1}{1}{1:14:42}\ts{1}{1}{1:15:04}\ts{1}{1}{1:15:25}\ts{1}{1}{1:15:46}. Schuller states that this can again be proven in a few lines, and does not for lack of time (already 3 minutes
over). The proof, \emph{not said in class}:
\begin{proposition}[Restrictions are continuous]\label{prop:c01-restriction}
Let $f\colon M\to N$ be continuous and $S\subseteq M$ carry the subset topology. Then $f|_S\colon S\to N$ is continuous.
\end{proposition}
\begin{proof}[Proof (not said in class)]
For $V\in\mathcal O_N$, $(f|_S)^{-1}(V)=\{s\in S\mid f(s)\in V\}=f^{-1}(V)\cap S$, and $f^{-1}(V)\in\mathcal O_M$ by continuity of $f$, so $f^{-1}(V)\cap S\in\mathcal O|_S$.
\end{proof}

\paragraph{Wrap-up.} A set was taken, and a topology established on it in order to talk about continuity \ts{1}{1}{1:15:46}\ts{1}{1}{1:16:06}. After the lecture the problem sheets for lectures 1 and
2 are on the tables (after the coffee break); one takes a sheet from each stack and comes to the tutorial in the afternoon in the bank building, lower floor: one is divided into groups
and works through a worksheet under the advice of the tutors; ``it will really make everything very clear'' and one is fully prepared for the next round of lectures on Wednesday.
There is half an hour coffee break, and the next lecture, at a quarter to, is on \emph{topological manifolds} \ts{1}{1}{1:16:06}\ts{1}{1}{1:16:26}\ts{1}{1}{1:16:47}.

\begin{history}[Where the definition comes from]
Schuller does not give the history. \emph{Checked.} Hausdorff's \emph{Grundz\"uge der Mengenlehre} (1914) introduced abstract topological spaces by axioms on
\emph{neighbourhoods}; the equivalent formulation by open sets, the one used here, became standard in the following decades (e.g.\ in Alexandroff and
Hopf's \emph{Topologie}, 1935). (Not said in class.)
\end{history}

% ==================================================================
\section{Tutorial 1}\label{sec:t01-1}

Tutorial 1, 0:32:35, \seg{1}{2}{video}: given by a tutor, not by Schuller. The sheet (Sheet 1, topology) is cited and paraphrased, not copied; see
\url{https://tales.mbivert.com/on-heraeus-winter-school-tutorials/} (and \url{https://tales.mbivert.com/heraeus/Sheet_01.pdf}). The tutorial recalls that the first lecture was about topology
and topological spaces and that the first exercise consists of multiple-choice questions, designed to spark discussion and as a self-test: tick the correct statements, not the incorrect
ones \ts{1}{2}{0:00}\ts{1}{2}{0:20}\ts{1}{2}{0:40}.

\subsection*{Exercise 1: multiple-choice questions}

\paragraph{(A) A topological space\dots}
\begin{itemize}
  \item \dots is defined by a set, a topology and an atlas: \textbf{false}. It is only a set equipped with a topology; an atlas is additional structure needed for a manifold, introduced later \ts{1}{2}{0:40}\ts{1}{2}{1:01}.
  \item \dots is a set without any further structure: \textbf{false}, the topology is further structure \ts{1}{2}{1:01}\ts{1}{2}{1:23}.
  \item \dots defines a notion of open sets: \textbf{true}, the topology says which sets are open \ts{1}{2}{1:23}\ts{1}{2}{1:44}.
  \item \dots always has an integer dimension: \textbf{false}; one cannot speak of dimension when dealing with a topological space \ts{1}{2}{1:44}.
  \item \dots allows one to check the continuity of a map from the underlying set to itself: \textbf{true}; a topological space is the minimal structure required to define continuity of a map
        \ts{1}{2}{2:05}.
\end{itemize}
(So two correct answers in part A, the third and the fifth \ts{1}{2}{2:25}.)

\paragraph{(B) The chaotic topology on a set $M$\dots}
\begin{itemize}
  \item \dots cannot be defined on the natural numbers: \textbf{false}, of course one can equip $\N$ with it \ts{1}{2}{2:25}\ts{1}{2}{2:48}.
  \item \dots consists of all subsets of $M$: \textbf{false}; that is the \emph{discrete} topology, the chaotic one consists only of $\emptyset$ and $M$ \ts{1}{2}{2:48}.
  \item \dots contains the empty set: \textbf{true}, this is the first axiom \ts{1}{2}{3:08}.
  \item \dots is the coarsest topology on $M$: \textbf{true}. Definition of ``coarser'' \ts{1}{2}{3:08}\ts{1}{2}{3:32}:
\end{itemize}
\begin{definition}[Coarser and finer topology (tutorial)]\label{def:c01-coarser}
A topology $\mathcal O_2$ on $M$ is \textbf{coarser} than a topology $\mathcal O_1$ (and $\mathcal O_1$ is \textbf{finer}) if $\mathcal O_2\subseteq\mathcal O_1$
\ts{1}{2}{3:32}\ts{1}{2}{3:57}.
\end{definition}
Being contained in another set is only a \emph{partial} order: there can be topologies that cannot be said to be subsets of each other, so one cannot compare them; only for some pairs can one say that one
is coarser. The chaotic topology $\{\emptyset,M\}$ is clearly the coarsest topology on $M$ \ts{1}{2}{3:57}\ts{1}{2}{4:18}\ts{1}{2}{4:38}\ts{1}{2}{4:59}.
\begin{itemize}
  \item \dots makes every map $f\colon N\to M$ continuous, where $N$ carries an arbitrary topology: \textbf{true}. The open sets of the target $M$ are $\emptyset$ and $M$; their pre-images are
        $\emptyset$ and $N$, both open in any topological space $N$ \ts{1}{2}{4:59}\ts{1}{2}{5:20}\ts{1}{2}{5:42}\ts{1}{2}{6:03}.
\end{itemize}

\paragraph{(C) A map $f\colon M\to N$ between topological spaces $(M,\mathcal O_M)$ and $(N,\mathcal O_N)$\dots}
\begin{itemize}
  \item \dots continuity can only be defined if $M=\R^m$, $N=\R^n$: \textbf{false}; arbitrary sets with topologies are allowed \ts{1}{2}{6:03}\ts{1}{2}{6:23}.
  \item \dots can for some maps be arranged to coincide with the undergraduate (epsilon--delta) notion: \textbf{true}; for $M=\R^m$, $N=\R^n$ with the standard topologies the topological notion
        equals the $\varepsilon$--$\delta$ one \ts{1}{2}{6:43}\ts{1}{2}{7:04}\ts{1}{2}{7:25}.
  \item \dots continuity is defined only for some topologies: \textbf{false}, it is defined for any topology \ts{1}{2}{7:25}\ts{1}{2}{7:46}.
  \item \dots continuity depends only on $\mathcal O_M$: \textbf{false}; pre-images of open sets of the target have to be open in the domain, so one needs both $\mathcal O_N$ and $\mathcal O_M$
        \ts{1}{2}{7:46}\ts{1}{2}{8:07}.
  \item \dots choosing the discrete topology on $M$ makes every map $M\to N$ continuous: \textbf{true}; all subsets of $M$ are open, so any pre-image is open, whatever topology $N$ carries
        \ts{1}{2}{8:28}\ts{1}{2}{8:49}\ts{1}{2}{9:10}.
\end{itemize}

\paragraph{(D) A subset $U\subseteq M$ of a topological space\dots}
\begin{itemize}
  \item \dots may be open and not open at the same time: \textbf{false} (nonsense; it is either open or not) \ts{1}{2}{9:10}\ts{1}{2}{9:30}.
  \item \dots may be open but not closed: \textbf{true}; example $M=\R$ with the standard topology, $U=(0,1)$ open; its complement $(-\infty,0]\cup[1,\infty)$ is not open
        \ts{1}{2}{9:30}\ts{1}{2}{9:52}\ts{1}{2}{10:13}\ts{1}{2}{10:34}. (The transcript has the answer ``this is wrong it is right''; the statement is true.)
  \item \dots may be closed but not open: \textbf{true}; $U=[0,1]$ in $\R$ \ts{1}{2}{10:34}\ts{1}{2}{10:55}.
  \item \dots may be open and closed: \textbf{true}; shown in the lecture: the set $M$ itself is open and closed in any topology \ts{1}{2}{11:16}\ts{1}{2}{11:37}.
  \item \dots may be neither open nor closed: \textbf{true}; $U=[0,1)$ in $\R$ \ts{1}{2}{11:37}\ts{1}{2}{11:59}.
\end{itemize}

\subsection*{Exercise 2: topologies on a simple set}

The first question is to write down the definition of a topology; this is \cref{def:c01-topology} (empty set and $M$ in $\mathcal O$; closed under pair-wise intersections; closed under
unions of arbitrary families indexed by some set) \ts{1}{2}{11:59}\ts{1}{2}{12:21}\ts{1}{2}{12:42}\ts{1}{2}{13:04}\ts{1}{2}{13:25}\ts{1}{2}{13:48}\ts{1}{2}{14:09}.

Take $M=\{1,2,3,4\}$. The sheet gives a first candidate $\mathcal O_1$ consisting of $\emptyset$, $\{1\}$ and $M$ (and a second one, $\mathcal O_2$, containing among others $\{1\}$ and $\{2\}$).
\begin{itemize}
  \item $\mathcal O_1=\{\emptyset,\{1\},M\}$ \emph{is} a topology: the first criterion holds; all intersections ($\emptyset\cap\{1\}=\emptyset$, $\{1\}\cap M=\{1\}$, \dots) are in $\mathcal O_1$;
        all unions again land in $\mathcal O_1$ \ts{1}{2}{14:09}\ts{1}{2}{14:31}\ts{1}{2}{14:52}\ts{1}{2}{15:13}\ts{1}{2}{15:36}\ts{1}{2}{15:57}\ts{1}{2}{16:18}\ts{1}{2}{16:39}.
  \item $\mathcal O_2$ is \emph{not} a topology: the first and second criteria hold, but the union of $\{1\}$ and $\{2\}$ is $\{1,2\}\notin\mathcal O_2$, violating the third
        \ts{1}{2}{16:39}\ts{1}{2}{17:01}\ts{1}{2}{17:22}\ts{1}{2}{17:42}.
  \item Are there other topologies on $M$? Definitely: the chaotic one, the discrete one and many more \ts{1}{2}{17:42}\ts{1}{2}{18:04}\ts{1}{2}{18:25}.
\end{itemize}

\subsection*{Exercise 3: continuous maps}

The definition of continuity: pre-images of open sets in the target are open in the domain \ts{1}{2}{18:46}\ts{1}{2}{19:11}. Let again $M=\{1,2,3,4\}$ and consider the identity map $\mathrm{id}_M$, $m\mapsto m$
\ts{1}{2}{19:34}.
\begin{itemize}
  \item Domain with the \emph{chaotic} topology, target with a topology $\mathcal O_{\mathrm{target}}=\{\emptyset,\{1\},M\}$: the pre-image of the open set $\{1\}$ is $\{1\}$, which is not in the chaotic
        topology, so $\mathrm{id}$ is \emph{not continuous} \ts{1}{2}{19:55}\ts{1}{2}{20:16}\ts{1}{2}{20:36}\ts{1}{2}{20:57}.
  \item The same set map in the opposite direction (the inverse identity), now with domain $\mathcal O_{\mathrm{target}}$ and target the chaotic topology: the open sets of the target are $\emptyset$ and $M$, whose
        pre-images are $\emptyset$ and $M$, both in the domain's topology. So the inverse identity \emph{is} continuous \ts{1}{2}{21:22}\ts{1}{2}{21:43}\ts{1}{2}{22:04}\ts{1}{2}{22:25}\ts{1}{2}{22:47}\ts{1}{2}{23:08}.
\end{itemize}
This again shows how the notion of continuity is directly affected by the chosen topologies on domain and target \ts{1}{2}{23:08}\ts{1}{2}{23:30}. (Compare \cref{ex:c01-f-g}.)

\subsection*{Exercise 4: the standard topology on \texorpdfstring{$\R^d$}{Rd}}

The last exercise is about the standard topology: sketch given real intervals and decide whether they are open \ts{1}{2}{23:30}\ts{1}{2}{23:50}.
\begin{itemize}
  \item $(0,1)$, endpoints excluded: open, since around every point there is an open ball inside \ts{1}{2}{23:50}\ts{1}{2}{24:11}.
  \item $[0,1)$: not open, at the point $0$ there is no open ball inside the interval \ts{1}{2}{24:11}\ts{1}{2}{24:32}\ts{1}{2}{24:53}.
  \item $(0,1]$: not open, now one fails at the point $1$ \ts{1}{2}{24:53}\ts{1}{2}{25:14}.
  \item $[0,1]$: not open \ts{1}{2}{25:14}\ts{1}{2}{25:35}.
  \item A union of two open intervals (in the sheet: an open interval together with a further open interval; the numbers are garbled in the captions [?]) is open: unions of open sets are open (axiom three), and
        again one finds balls around every point \ts{1}{2}{25:35}\ts{1}{2}{25:57}\ts{1}{2}{26:18}.
\end{itemize}
Then six subsets of $\R^2$ with the standard topology (the frame drawn around them has no meaning, it stands for all of $\R^2$) \ts{1}{2}{26:18}\ts{1}{2}{26:39}\ts{1}{2}{27:00}:
\begin{enumerate}
  \item open: around every point one finds a ball inside \ts{1}{2}{27:00}\ts{1}{2}{27:21};
  \item open, likewise \ts{1}{2}{27:21}\ts{1}{2}{27:41};
  \item not open: there is a boundary point around which every ball sticks out of the set \ts{1}{2}{27:41}\ts{1}{2}{28:02};
  \item an open interval of $\R$ placed inside $\R^2$ (a segment of a line): \emph{not} open in $\R^2$, no ball of $\R^2$ lies completely in it, although the interval is open in $\R$ \ts{1}{2}{28:02}\ts{1}{2}{28:24}\ts{1}{2}{28:45};
  \item the complement of the second set: it is closed, but not open (for some points there is no ball) \ts{1}{2}{28:45}\ts{1}{2}{29:07}\ts{1}{2}{29:28};
  \item again a union of open sets, hence open \ts{1}{2}{29:28}\ts{1}{2}{29:49}.
\end{enumerate}
Finally a function $f\colon\R\to\R$ given by a graph with a jump, both sides carrying the standard topology: is it continuous? Take an open set in the target; its pre-image
is an interval that includes one of its endpoints (the exact numbers are garbled in the captions [?]); by the previous question such an interval is not open in the standard topology, so $f$ is
\emph{not continuous}, as expected from undergraduate intuition since the function makes a jump. For maps on the reals with the standard topology, the topological notion of continuity coincides with
the undergraduate one \ts{1}{2}{29:49}\ts{1}{2}{30:10}\ts{1}{2}{30:31}\ts{1}{2}{30:54}\ts{1}{2}{31:20}\ts{1}{2}{31:42}\ts{1}{2}{32:03}\ts{1}{2}{32:24}.

\begin{coursenote}[Tutorial sheet]
The solutions above follow the tutor's order of the sheet. The sheet itself (Sheet 1) is not reproduced here: copyright; see the link at the top of this section.
\end{coursenote}

% ==================================================================
\section*{Key concepts}
\begin{keyconcepts}
\textbf{Retrieval questions.}
\begin{enumerate}
  \item Why does one need more than a set to talk about continuity of curves? What is the weakest structure that suffices?
  \item State the three axioms of a topology. Which of them fail for $\{\emptyset,\{1\},\{2\},\{1,2,3\}\}$ on $\{1,2,3\}$?
  \item What are the chaotic and the discrete topologies? Why are they useful although ``useless''?
  \item Define soft balls and the standard topology on $\R^d$. Why is it not special?
  \item Define open and closed sets. Give a set that is open and closed, and one that is neither.
  \item Define continuity via pre-images. Does it depend on the topologies of both domain and target?
  \item Give a continuous bijection whose inverse is not continuous. What is a homeomorphism?
  \item Prove that the composition of continuous maps is continuous.
  \item What is the subset topology, and why is it the natural one for restrictions of maps?
\end{enumerate}
\textbf{Concepts.}
\begin{itemize}
  \item Topology $\mathcal O\subseteq\mathcal P(M)$ with $\emptyset,M\in\mathcal O$, finite intersections, arbitrary unions (\cref{def:c01-topology}); topological space $(M,\mathcal O)$ (\cref{def:c01-open-closed}).
  \item Chaotic $\{\emptyset,M\}$ and discrete $\mathcal P(M)$ topology (\cref{def:c01-extreme}); coarser and finer (\cref{def:c01-coarser}).
  \item $B_r(p)$ and $\mathcal O_{\mathrm{standard}}$: $U$ open iff each point has a ball inside (\cref{def:c01-standard,prop:c01-std-ok}).
  \item Open sets are the elements of $\mathcal O$; closed sets have open complements; the two notions are independent (\cref{rem:c01-warning}).
  \item $f$ continuous iff pre-images of open sets are open (\cref{def:c01-continuous}); homeomorphism (\cref{def:c01-homeo}, \cref{ex:c01-f-g}).
  \item $g\circ f$ continuous if $f,g$ are (\cref{thm:c01-composition}).
  \item Subset topology $\mathcal O|_S=\{U\cap S\}$ and continuity of restrictions (\cref{def:c01-subset,prop:c01-subset-ok,prop:c01-restriction}).
\end{itemize}
\end{keyconcepts}

\flushside
\end{paracol}
